\documentclass[11pt]{article}

\usepackage[margin=1in]{geometry}
\usepackage[T1]{fontenc}
\usepackage[utf8]{inputenc}
\usepackage{microtype}
\usepackage{booktabs}
\usepackage{longtable}
\usepackage{array}
\usepackage{hyperref}
\usepackage{tikz}
\usepackage[numbers,sort&compress]{natbib}

\title{Workplace Determinants of Metabolic Syndrome in Working Adults: Systematic Review}
\author{Gideon Ofosu Kyere\\
Kwame Nkrumah University of Science and Technology (KNUST), Kumasi, Ghana\\
ORCID: 0009-0003-9848-8437}
\date{}

\begin{document}
\maketitle

\begin{abstract}
\textbf{Objectives:}
To synthesise evidence on workplace stress, diet, physical activity, sedentary behaviour and metabolic syndrome (MetS) among working adults, and to evaluate workplace-based lifestyle interventions targeting MetS or its established components.

\textbf{Methods:}
Peer-reviewed studies published from 2015 to 2026 were identified using a workplace-focused search. Eligible studies examined employed adults or occupational populations and evaluated relevant exposures or workplace lifestyle interventions in relation to MetS or established components. Randomised trials, cohort studies, analytical cross-sectional studies and quasi-experimental interventions were eligible. Design-specific Joanna Briggs Institute tools were used for critical appraisal. Heterogeneous evidence was synthesised narratively.

\textbf{Results:}
Prospective evidence linked adverse or increasing work stress with greater MetS risk, while cross-sectional findings were more mixed. Higher leisure-time physical activity was associated with lower MetS incidence or prevalence, whereas prolonged sedentary exposure was associated with greater risk. Healthier dietary patterns were generally associated with more favourable metabolic profiles. Workplace interventions involving exercise, dietary modification, counselling or multicomponent lifestyle support improved waist circumference, blood pressure, glucose regulation, lipid profiles or MetS severity in several studies.

\textbf{Conclusions:}
Workplace stress, diet, physical activity and sedentary behaviour are potentially modifiable correlates of MetS among working adults. Workplace lifestyle interventions can improve metabolic outcomes, although confidence is greatest for randomised and prospective evidence.
\end{abstract}

\section*{What is already known on this topic}
Metabolic syndrome is common in working populations, and occupational stress, diet, physical inactivity and sedentary behaviour have each been studied as potential contributors. However, the evidence is dispersed across different occupational groups and study designs.

\section*{What this study adds}
This systematic review integrates prospective, cross-sectional and workplace-intervention evidence across occupational stress, diet/nutrition, physical activity/sedentary behaviour and multicomponent workplace programmes. Prospective studies support an association between adverse work stress and metabolic syndrome risk, while healthier diet and greater physical activity are generally associated with more favourable metabolic profiles; workplace interventions can improve several metabolic-syndrome components.

\section*{How this study might affect research, practice or policy}
Workplace prevention strategies may be strengthened by combining psychosocial risk management, healthier food environments, opportunities for physical activity, reduced prolonged sitting and targeted cardiometabolic screening. Future occupational-health research should prioritise consistent metabolic-syndrome definitions and stronger longitudinal and randomised designs.

% ----------------------------------------------------------------
% Scientific body
% Copy/paste the frozen body from manuscript/main.tex beginning at
% \section{Introduction} and ending immediately before the current
% declarations/bibliography block. Do not edit substantive results.
% ----------------------------------------------------------------

\section*{Registration and protocol}
This review update was not prospectively registered. No prospectively registered protocol was available; however, the revised eligibility criteria, search framework, screening rules, extraction structure and appraisal workflow were defined before the final evidence synthesis and are archived in the project repository.

\section*{Reporting bias and certainty of evidence}
Formal assessment of publication bias was not undertaken because no pooled meta-analysis was performed and the evidence was highly heterogeneous across study designs, occupational populations, exposures, interventions and outcomes. A formal GRADE certainty-of-evidence framework was not applied. Confidence in the findings was interpreted in relation to study design, directness, consistency across studies and design-specific Joanna Briggs Institute appraisal findings.

\section*{Funding}
No external funding was received for this systematic review.

\section*{Competing interests}
None declared.

\section*{Ethics approval}
Ethics approval was not required because this systematic review synthesises previously published studies and did not involve collection of new participant-level data.

\section*{Data availability}
Study-level extraction tables, appraisal files, search documentation and reproducibility materials are available in the project repository:
\url{https://github.com/OG-Kyere/workplace-stress-metabolic-syndrome-systematic-review}.

\section*{Author contributions}
Gideon Ofosu Kyere conceived the review update, developed the revised workplace-focused eligibility framework, conducted the updated literature search and screening, extracted and synthesised the evidence, performed methodological appraisal, and drafted and revised the manuscript.

\section*{Use of artificial intelligence}
Generative artificial-intelligence tools were used for editorial and organisational assistance during manuscript development. All eligibility decisions, methodological judgments, extracted study results, interpretations and final manuscript content were reviewed by the author against the underlying sources.

\bibliographystyle{unsrtnat}
\bibliography{references_verified}

\end{document}
